\documentclass[12pt,a4paper]{article}

\usepackage{amsmath,amssymb}
\usepackage{graphicx}
\usepackage{geometry}
\usepackage{hyperref}
\usepackage{float}
\usepackage{natbib}

\geometry{margin=1in}

\title{005\_BIMOD: A Discrete Four-Locus Polarity Classification Methodology}
\author{
DuckJin Yoon (Solvaram) \\
\textit{Eight-Constitutional Research Platform} \\
\texttt{with methodological clarification by DeepSeek AI}
}
\date{Version 3.2 -- February 24, 2026}

\begin{document}

\maketitle

\begin{abstract}
This paper formalizes the discrete methodological foundation of the BIMOD framework. 
A nasion-centered reference structure defines four anatomical loci. 
Each locus admits a binary cranial polarity state (+1 or −1), yielding an eight-state classification system. 
The model is presented as a discrete structural taxonomy independent of higher-level physical hypotheses.
\end{abstract}

\section{Methodological Scope}

The 005\_BIMOD paper defines only the classification architecture. 
It does not assert a continuous vector field, magnetic mechanism, or biophysical resonance model. 
These higher-level interpretations are addressed in BIMOD-500, which integrates physical hypotheses with the taxonomic base established here.

\section{Structural Definition}

Let the nasion be the anatomical reference origin, a craniofacial landmark known for its positional stability across individuals \citep{nasion_wikipedia}.

Define four loci symmetrically positioned around the nasion:

\[
L = \{L_1, L_2, L_3, L_4\}
\]

Each locus admits a binary polarity state:

\[
P(L_j) \in \{+1,-1\}
\]

The classification space is defined as the Cartesian product:

\[
\mathcal{C} = L \times \{+1,-1\}
\]

Thus the total number of classification states:

\[
|\mathcal{C}| = 4 \times 2 = 8
\]

This discrete structure ensures categorical exclusivity: each individual corresponds to exactly one of the eight possible states.

\section{Geometric Reference}

\begin{figure}[H]
\centering
\includegraphics[width=0.7\textwidth]{005_NasionAI.jpg}
\caption{
Full-face nasion reference image showing the approximate locations of the four measurement loci. 
Each locus (L1–L4) admits two polarity states, forming an eight-state classification system. 
The actual detection procedure follows BIMOD protocol rather than visual inspection.
}
\label{fig:nasion}
\end{figure}

\section{Label Layer and Optional Mappings}

The eight states may optionally be mapped to constitutional labels, such as:

\begin{itemize}
    \item 1S (Hepatonia), 1N (Pulmotonia)
    \item 2N (Cholecystonia), 2S (Colonotonia)
    \item 3N (Gastrotonia), 3S (Vesicotonia)
    \item 4N (Pancreotonia), 4S (Renotonia)
\end{itemize}

However, such interpretive mappings belong to higher-level theoretical layers (Track B in BIMOD-500) and are not intrinsic to this methodological definition. The taxonomic structure stands independently of any particular interpretation.

\section{Statistical Reproducibility}

Classification reproducibility may be evaluated using a binomial framework \citep{binomialtest}:

\[
P(X \ge k | N, p=0.5) = \sum_{i=k}^{N} \binom{N}{i} (0.5)^N
\]

For \(N=20\) trials, the threshold for significance at \(\alpha=0.05\) (one-sided) is \(k=15\), yielding \(P(X \ge 15) = 0.0207\).

Empirical validation requires:
\begin{itemize}
    \item Blinded protocols (practitioner unaware of sample polarity)
    \item RAW image preservation (no AI enhancement or filtering)
    \item Independent replication across multiple practitioners
\end{itemize}

\section{Relation to Physical Hypotheses}

Biogenic magnetic nanoparticles have been reported in biological systems \citep{Kirschvink1992}, and magnetoreception models based on geomagnetic interaction have been proposed \citep{Ritz2000}. However, the taxonomic system defined here does not depend on any specific physical mechanism. It serves as a structural framework that can accommodate various explanatory models.

\section{Boundary Condition}

This paper defines a discrete taxonomy only. 
Physical explanatory models, continuous field formalizations, and empirical applications are addressed in BIMOD-500 and other papers in the series.

\section{Conclusion}

The nasion-based eightfold polarity system is presented as a discrete methodological classification framework, open to theoretical interpretation and empirical examination while remaining independent of both.

\begin{thebibliography}{99}

\bibitem{Kirschvink1992}
Kirschvink, J. L. et al. (1992). Magnetite in human tissues: A mechanism for biomagnetic detection? \textit{Bioelectromagnetics}.

\bibitem{Ritz2000}
Ritz, T. et al. (2000). A model for photoreceptor-based magnetoreception in birds. \textit{Biophysical Journal}, 78(2), 707–718.

\bibitem{binomialtest}
Clopper, C. J. \& Pearson, E. S. (1934). The use of confidence or fiducial limits illustrated in the case of the binomial. \textit{Biometrika}, 26, 404–413.

\bibitem{nasion_wikipedia}
Nasion. Wikipedia. Accessed 2026.

\end{thebibliography}

\end{document}